%%%%%%%%%%%%%%%%%%%%%%%%%%%%%%%%%%%%%%%%%%%%%%%%%%%%%%%%%%%%
\section{性能}
\label{sec:performance}

我们通过使用 ECDSA-256 签名和 SHA-256 哈希在通用硬件（8 核、16GB RAM）上进行微基准测试（microbenchmark）来验证实际开销。\textbf{密码学操作}通常开销低于 $\sim$0.2ms（ECDSA 签名生成 $\sim$102$\mu$s，验证 $\sim$89$\mu$s，每次操作哈希链更新 $<$15$\mu$s）。\textbf{策略操作}：检索延迟取决于后端：内存（$\sim$5-50$\mu$s）、文件系统（$\sim$30-200$\mu$s）、Redis（$\sim$50-300$\mu$s）；策略交集取决于组合方式——典型 3-5 层深度的层级结构下保持在 100$\mu$s 以内。\textbf{自定义验证器}：确定性检查（PII 检测、路径验证）增加 30-400$\mu$s；基于 LLM 的验证器（提示安全性）增加 150-500ms，启用时占主导地位，但提供与密码学保证正交的语义分析。\textbf{端到端影响}：对于网络绑定操作（LLM 推理、远程 API），密码学开销（$\sim$0.2ms）相对于网络延迟（50-500ms）可以忽略不计。对于本地工具调用，不含 LLM 验证器时，PEP 验证总开销保持在 1ms 以下。审计日志通过缓冲 I/O 异步（asynchronous）进行。

配套论文提供了全面的系统评估，包括生产工作负载特征化、跨越全部九个框架的框架集成成本、大规模部署经验以及基于生产级基准测试的详细性能分析。
